%==============================================================
%  Figura 8.2 - Campo di applicazione della EN ISO 13849 per i
%  sistemi pneumatici
%  Adattamento in italiano della Figura 8.2 dell'IFA Report 2/2017
%  (stesso stile della Figura 8.1)
%
%  I simboli dei componenti riproducono quelli dell'originale
%  (ISO 1219): le coordinate sono ricavate dal disegno di partenza
%  (1 cm = 70 px, origine nell'angolo in alto a sinistra della
%  zona tratteggiata). Nessun cambio di corpo del font: vedere la
%  nota nella Figura 8.1.
%==============================================================
\begin{tikzpicture}[
  font=\normalsize,
  linea/.style={draw=black, line width=0.9pt},
  simbolo/.style={draw=black, line width=0.9pt},
  freccia/.style={simbolo, -{Latex[length=5pt,width=4pt]}},
  doppia/.style={simbolo, {Latex[length=5pt,width=4pt]}-{Latex[length=5pt,width=4pt]}},
  zona/.style={draw=black, line width=0.9pt, dash pattern=on 7pt off 5pt},
  pilota/.style={draw=black, line width=0.7pt, dash pattern=on 3pt off 3pt},
  testo/.style={align=center, inner sep=1pt},
  nodo/.style={circle, fill=black, inner sep=0pt, minimum size=4pt}
]

%--------------------------------------------------------------
% Zone tratteggiate: componenti della funzione di sicurezza,
% componenti contro le fluttuazioni dell'energia, unita' di
% manutenzione
%--------------------------------------------------------------
\draw[zona] (0,0) rectangle (11.09,-16.71);
\draw[zona] (0,-5.89) -- (11.09,-5.89);
\draw[zona] (0,-9.60) -- (11.09,-9.60);
\draw[zona] (1.87,-10.86) rectangle (10.54,-14.83);

\node[testo, text width=5.6cm] at (5.51,-2.86)
  {Componenti che realizzano\\ la funzione di sicurezza,\\ ad es. valvole};
\node[testo, text width=5.6cm] at (5.51,-7.83)
  {Componenti per la prevenzione\\ dei rischi in caso di\\ fluttuazioni dell'energia};
\node[anchor=east] at (1.72,-10.98) {0Z};
\node[anchor=east] at (10.90,-16.00) {\textquotedblleft Unit\`a di manutenzione\textquotedblright};

%--------------------------------------------------------------
% Cilindro 1A (elemento di azionamento): camicia, pistone e stelo
% a doppia linea, blocco di attacco con le due bocche
%--------------------------------------------------------------
\draw[simbolo] (4.44,2.17) rectangle (6.57,3.33);
\draw[simbolo] (4.67,2.17) -- (4.67,3.33);
\draw[simbolo] (5.14,2.17) -- (5.14,3.33);
\draw[simbolo, fill=white] (5.14,2.61) rectangle (7.29,2.86);
\draw[linea] (4.50,2.17) -- (4.50,1.31) -- (5.29,1.31) -- (5.29,0);
\draw[linea] (6.50,2.17) -- (6.50,1.31) -- (5.71,1.31) -- (5.71,0);

\node[anchor=east] at (4.30,3.31) {1A};
\node[anchor=east] at (4.30,2.14) {Elementi di azionamento};

%--------------------------------------------------------------
% Linea di alimentazione pneumatica (interrotta dal filtro 0Z10 e
% dalla valvola 0V10)
%--------------------------------------------------------------
\draw[linea] (0.50,-9.60) -- (0.50,-13.43) -- (6.29,-13.43);
\draw[linea] (7.60,-13.43) -- (8.57,-13.43);
\draw[linea] (9.54,-13.43) -- (11.29,-13.43);
% sorgente di aria compressa
\draw[simbolo] (11.29,-13.43) -- (11.69,-13.20) -- (11.69,-13.66) -- cycle;

%--------------------------------------------------------------
% 0Z11: manometro
%--------------------------------------------------------------
\draw[simbolo] (2.96,-12.49) circle (0.46);
\draw[freccia] (3.15,-12.70) -- (2.74,-12.30);
\draw[linea] (2.96,-12.95) -- (2.96,-13.43);
\node[nodo] at (2.96,-13.43) {};
\node[anchor=south] at (2.96,-11.86) {0Z11};

%--------------------------------------------------------------
% 0V11: regolatore di pressione con molla regolabile, sfiato
% secondario e pilotaggio
%--------------------------------------------------------------
\draw[simbolo] (4.33,-12.73) rectangle (5.27,-13.66);
\draw[doppia] (4.37,-13.43) -- (5.23,-13.43);
% molla verticale sopra la valvola
\draw[simbolo] (5.07,-12.73) -- (5.50,-12.64) -- (4.93,-12.54) -- (5.50,-12.44)
              -- (4.93,-12.34) -- (5.50,-12.24) -- (5.11,-12.16);
% freccia di regolazione
\draw[freccia] (5.53,-12.56) -- (4.46,-12.30);
% sfiato secondario
\draw[simbolo] (5.27,-12.86) -- (5.50,-12.98) -- (5.27,-13.10);
% pilotaggio
\draw[pilota] (4.33,-13.49) -- (4.11,-13.66) -- (4.11,-14.13) -- (5.04,-14.13)
             -- (5.04,-13.66);
\node[anchor=south] at (4.83,-11.86) {0V11};

%--------------------------------------------------------------
% 0Z10: filtro con separatore d'acqua e indicatore di intasamento
%--------------------------------------------------------------
\draw[simbolo] (6.93,-12.79) -- (7.60,-13.43) -- (6.93,-14.09) -- (6.29,-13.43) -- cycle;
\draw[pilota] (6.93,-12.79) -- (6.93,-13.66);
\draw[simbolo] (6.54,-13.66) -- (7.33,-13.66) -- (6.93,-13.96) -- cycle;
\draw[linea] (6.93,-14.09) -- (6.93,-15.29);
% indicatore di intasamento
\draw[simbolo] (6.34,-12.06) rectangle (7.53,-12.77);
\draw[freccia] (6.88,-12.60) -- (6.49,-12.21);
\draw[simbolo] (7.19,-12.39) circle (0.16);
\draw[simbolo] (7.08,-12.28) -- (7.30,-12.50);
\draw[simbolo] (7.08,-12.50) -- (7.30,-12.28);
\node[anchor=south] at (6.93,-11.86) {0Z10};

%--------------------------------------------------------------
% 0V10: valvola di intercettazione manuale 3/2 con scarico
%--------------------------------------------------------------
\draw[simbolo] (8.57,-12.24) rectangle (9.54,-14.14);
\draw[simbolo] (8.57,-13.20) -- (9.54,-13.20);
% posizione di passaggio (porta chiusa a destra)
\draw[doppia] (8.61,-12.50) -- (9.49,-12.50);
\draw[simbolo] (9.24,-12.89) -- (9.24,-13.03);
\draw[simbolo] (9.24,-12.96) -- (9.54,-12.96);
% posizione di scarico (porta chiusa sulla linea)
\draw[doppia] (8.60,-13.43) -- (9.54,-13.89);
\draw[simbolo] (9.24,-13.36) -- (9.24,-13.51);
\draw[simbolo] (9.24,-13.43) -- (9.54,-13.43);
\draw[linea] (9.54,-13.89) -- (9.97,-13.89);
\draw[simbolo] (9.97,-13.79) -- (9.97,-14.00) -- (10.19,-13.89) -- cycle;
% azionamento manuale
\draw[linea] (9.24,-12.24) -- (9.24,-11.31);
\draw[linea] (9.00,-12.24) -- (9.00,-11.91) -- (9.11,-11.79) -- (9.00,-11.67)
            -- (9.00,-11.31);
\draw[simbolo] (8.89,-11.79) -- (8.97,-11.79);
\draw[simbolo] (8.89,-11.31) -- (9.34,-11.31);
\node[anchor=north] at (9.06,-14.26) {0V10};

%--------------------------------------------------------------
% Graffe e campi di applicazione
%--------------------------------------------------------------
\draw[linea, decorate, decoration={brace, amplitude=8pt}]
  (13.11,-5.64) -- (13.11,-0.29);
\node[testo, text width=5.2cm, anchor=west] at (13.77,-2.96)
  {Campo di applicazione della\\ parte del sistema di comando\\ legata alla sicurezza};

\draw[linea, decorate, decoration={brace, amplitude=8pt}]
  (13.11,-16.43) -- (13.11,-6.14);
\node[testo, text width=5.2cm, anchor=west] at (13.77,-11.46)
  {Possibile rilevanza per il\\ rispetto dei principi di\\ sicurezza di base e collaudati};

%--------------------------------------------------------------
% Cornice
%--------------------------------------------------------------
\draw[draw=black, line width=1pt]
  ([shift={(-0.45,0.45)}]current bounding box.north west) rectangle
  ([shift={(0.45,-0.45)}]current bounding box.south east);

\end{tikzpicture}
